\subsection{Extension of the base field in regular algebras (nonzero characteristic)}

Let \(k\) be a ring and \(\rho: A \to B\) a homomorphism of \(k\)-algebras. The homomorphism \(\rho\) induces an A-linear map \(\Omega(\rho): \Omega_k(A) \to \Omega_k(B)\), and consequently a B-linear map \(\Omega_0(\rho): B \otimes_A \Omega_k(A) \to \Omega_k(B)\) (A, III, p.~135). Let \(\mathbf{T} = (\mathrm{T}_i)_{i \in I}\) be a family of indeterminates, and \(\mathbf{t} = (t_i)_{i \in I}\) a family of elements of B; for every polynomial \(f = \sum_{\alpha \in \mathbf{N}^{(I)}} c_\alpha \mathbf{T}^\alpha\) of A{[} \(\mathbf{T}\) {]}, let us denote by \(d^{A} f(\mathbf{t})\) the element \(\sum_{\alpha} \mathbf{t}^\alpha \otimes dc_\alpha\) of \(B \otimes_A \Omega_k(A)\).

Lemma 1.—Suppose that the A-algebra B admits a generating family \(\mathbf{t} = (t_i)_{i\in I}\), subject to the relations \(f_\lambda \in \mathrm{A}[\mathbf{T}]\) (\(\lambda \in \Lambda\)). The B-linear homomorphism

\[
\psi : (\mathrm{B} \otimes_ {\mathrm{A}} \Omega_ {k} (\mathrm{A})) \oplus \mathrm{B} ^ {(\mathrm{I})} \longrightarrow \Omega_ {k} (\mathrm{B})
\]

defined by \(\psi(\alpha,(b_i)) = \Omega_0(\rho)(\alpha) + \sum_{i\in I}b_i\,dt_i\) is surjective; its kernel is generated by the elements \(n_{\lambda} = \left(d^{A}f_{\lambda}(\mathbf{t}),\left(\frac{\partial f_{\lambda}}{\partial\mathrm{T}_{i}}(\mathbf{t})\right)_{i\in I}\right)\) for \(\lambda \in \Lambda\).

Consider the sequence of B-modules and B-linear maps

\[
\mathrm{B} ^ {(\Lambda)} \xrightarrow {\varphi} (\mathrm{B} \otimes_ {\mathrm{A}} \Omega_ {k} (\mathrm{A})) \oplus \mathrm{B} ^ {(\mathrm{I})} \xrightarrow {\psi} \Omega_ {k} (\mathrm{B}) \longrightarrow 0,
\]

where \(\varphi\) is the homomorphism such that \(\varphi(e_{\lambda}) = n_{\lambda}\) ; it remains to show that this sequence is exact. By A, II, p.~36, th. 1, it suffices to prove that, for every B-module M, the sequence

\[
0 \to \operatorname{Hom} _ {\mathrm{B}} (\Omega _{k}(\mathrm{B}), \mathrm{M}) \xrightarrow {\operatorname{Hom}(\psi,1)} \operatorname{Hom} _{\mathrm{B}} ((\mathrm{B} \otimes_{\mathrm{A}} \Omega _{k}(\mathrm{A})) \oplus \mathrm{B}^{(\mathrm{I})}, \mathrm{M}) \xrightarrow {\operatorname{Hom}(\varphi,1)} \operatorname{Hom}_{\mathrm{B}}(\mathrm{B}^{(\Lambda)}, \mathrm{M})
\]

is exact. Given the universal property of the module of differentials (A, III, p.~134), this sequence identifies with

\[
0 \to \mathrm{D} _ {k} (\mathrm{B}, \mathrm{M}) \xrightarrow {\psi^ {\prime}} \mathrm{D} _ {k} (\mathrm{A}, \mathrm{M}) \oplus \mathrm{M} ^ {\mathrm{I}} \xrightarrow {\varphi^ {\prime}} \mathrm{M} ^ {\Lambda}
\]

where \(\psi^{\prime}(\mathrm{D}) = (\mathrm{D}\circ \rho,(\mathrm{D}(t_{i}))_{i\in I})\) and \(\varphi^{\prime}(\Delta,(m_{i})) = \big(f_{\lambda}^{\Delta}(\mathbf{t}) + \sum_{i}\frac{\partial f_{\lambda}}{\partial\mathrm{T}_{i}}(\mathbf{t})m_{i}\big)_{\lambda \in \Lambda}\) (in accordance with A, V, p.~121, for every polynomial \(f = \sum_{\alpha \in \mathbf{N}^{(I)}}c_{\alpha}\mathbf{T}^{\alpha}\) of A{[}T{]}, we denote \(f^{\Delta}(\mathbf{t})\) as the element \(\sum_{\alpha}\mathbf{t}^{\alpha}\Delta(c_{\alpha})\). Now the exactness of this sequence follows from loc. cit., prop. 1, taking into account that a derivation D : B → M is \(k\)-linear if and only if \(D\circ\rho\) is so.

Let A be a ring. There exists a unique structure of \(\mathbf{Z}\)-algebra on A; we simply denote \(\Omega(A)\) the A-module \(\Omega_{\mathbf{Z}}(\mathrm{A})\). If \(\rho : k \to \mathrm{A}\) is a ring homomorphism, we have a canonical exact sequence of A-modules (A, III, p.~136, prop. 21) \(\mathrm{A} \otimes_k \Omega(k) \to \Omega(\mathrm{A}) \to \Omega_k(\mathrm{A}) \to 0\).

Suppose that A contains a subfield, and let P be the prime subfield of A; then \(\Omega(\mathrm{P})\) is zero and the canonical homomorphism of A-modules \(\Omega(\mathrm{A}) \to \Omega_{\mathrm{P}}(\mathrm{A})\) is bijective. If moreover A has characteristic \(p \neq 0\) (which means by definition that \(p\) is a prime number, \(p1_{\mathrm{A}} = 0\), and \(1_{\mathrm{A}} \neq 0\)), then P is identified with \(\mathbf{F}_p\). Moreover, every derivation of A vanishes on the subring \(\mathbf{A}^p\); for every subring \(k\) of A contained in \(\mathbf{A}^p\) (and, in particular, for every perfect subfield \(k\) of A), the canonical map \(\Omega(\mathrm{A}) \to \Omega_k(\mathrm{A})\) is bijective.

Let A be a ring of characteristic \(p \neq 0\) and \((f_i)_{1 \leqslant i \leqslant n}\) a finite sequence of elements of A. Denote by \(A_n\) the quotient ring of the polynomial ring \(A[T_1, \ldots, T_n]\) by the ideal generated by the polynomials \(T_i^p - f_i\) , for \(1 \leqslant i \leqslant n\) .

Lemma 2.—Suppose the ring A is local and Noetherian. Then \(A_{n}\) is local and Noetherian. The following conditions are equivalent:

\begin{enumerate}
\def\labelenumi{(\roman{enumi})}
\item
  \(A_{n}\) is regular;
\item
  A is regular and the elements \(1 \otimes df_{i}\) of the \(\kappa_{A}\) -vector space \(\kappa_{A} \otimes_{A} \Omega(A)\) are linearly independent.
\end{enumerate}

The ring \(A_{n}\) is Noetherian (III, § 2, cor. 3 of th. 2). The A-module \(A_{n}\) is free, hence faithfully flat; if \(A_{n}\) is regular, then A is regular (§ 4, no. 5, prop. 8 b)). We shall reason by induction on n, the lemma being evident if n = 0.

\begin{enumerate}
\def\labelenumi{\Alph{enumi})}
\tightlist
\item
  Let us first treat the case n = 1, setting \(T_{1} = T\), \(f_{1} = f\). Let a denote the class of f in \(\kappa_{A}\) and distinguish two cases, according as a belongs or does not belong to \(\kappa_{A}^{p}\). If \(a \notin \kappa_{A}^{p}\), then the polynomial \(T^{p} - a\) is irreducible in \(\kappa_{A}\) (A, V, p.~24, lemma 1) and \(\kappa_{A} \otimes_{A} A_{1}\) is isomorphic to the field \(\kappa_{A}[T]/(T^{p} - a)\). The ideal \(m_{A}A_{1}\) of \(A_{1}\) is therefore maximal, so that the ring \(A_{1}\) is local (V, § 2, no. 1, prop. 1). If A is regular, \(A_{1}\) is regular (VIII, § 5, no. 1, prop. 1). According to A, V, p.~99, prop. 6, the element da of \(\Omega(\kappa_{A})\) is not zero; since it is the image under the canonical map \(\kappa_{A} \otimes_{A} \Omega(A) \to \Omega(\kappa_{A})\) of \(1 \otimes df\), the latter is not zero. This proves the lemma in this case.
\end{enumerate}

Now suppose that \(a\) belongs to \(\kappa_{\mathrm{A}}^{p}\). There therefore exists an element \(g\) of A such that \(f - g^{p} \in \mathfrak{m}_{\mathrm{A}}\). Set \(h = f - g^{p}\). Since \(\mathrm{T}^{p} - f = (\mathrm{T} - g)^{p} - h\), the A-algebra \(A_{1}\) is isomorphic to \(\mathrm{A}[\mathrm{T}] / (\mathrm{T}^{p} - h)\). According to VIII, § 5, no. 4, prop. 4, the ring \(A_{1}\) is local and, for it to be regular, it is necessary and sufficient that A be regular and that \(h\) not belong to \(\mathfrak{m}_{\mathrm{A}}^{2}\). Now, since \(\kappa_{A}\) is formally smooth over the prime field (no. 3, th. 1), the canonical map

\[
\bar {d}: \mathfrak {m} _ {\mathrm{A}} / \mathfrak {m} _ {\mathrm{A}} ^ {2} \to \kappa_ {\mathrm{A}} \otimes_ {\mathrm{A}} \Omega (\mathrm{A})
\]

is injective (no. 2, remark 1); but the image under \(\bar{d}\) of the class of \(h\) modulo \(\mathfrak{m}_{\mathrm{A}}^{2}\) is equal to \(1 \otimes dh = 1 \otimes d(f - g^{p}) = 1 \otimes df\). This proves the lemma in this second case and completes the proof of the case \(n = 1\).

\begin{enumerate}
\def\labelenumi{\Alph{enumi})}
\setcounter{enumi}{1}
\tightlist
\item
  Suppose \(n > 1\) . The ring \(\mathrm{A}_{1}\) is local and Noetherian by the case already treated. The \(\mathrm{A}_{1}\) -algebra \(\mathrm{A}_{n}\) is identified with the quotient of \(\mathrm{A}_{1}[\mathrm{T}_{2},\ldots ,\mathrm{T}_{n}]\) by the ideal generated by the \(\mathrm{T}_i^p -f_i\) , \(i\geqslant 2\) ; by the induction hypothesis, it is a local ring and condition (i) is equivalent to the conjunction of the following two conditions:
\end{enumerate}

(i') A \(_{1}\) is regular;

(i'\,') the elements \(1 \otimes df_2, \ldots, 1 \otimes df_n\) of the \(\kappa_{\mathrm{A}_1}\) -vector space \(\kappa_{\mathrm{A}_1} \otimes_{\mathrm{A}_1} \Omega(\mathrm{A}_1)\) are linearly independent.

But (i') is equivalent, as we have just seen, to

(ii') A is regular and the element \(1 \otimes df_{1}\) of the \(\kappa_{A}\)-vector space \(\kappa_{A} \otimes_{A} \Omega(A)\) is not zero.

By Lemma 1, the canonical homomorphism \(\mathrm{A}_{1} \otimes_{\mathrm{A}} \Omega(\mathrm{A}) \to \Omega(\mathrm{A}_{1})\) induces an isomorphism of \(((\mathrm{A}_{1} \otimes_{\mathrm{A}} \Omega(\mathrm{A})) / \mathrm{A}_{1}(1 \otimes df_{1})) \oplus \mathrm{A}_{1}\) onto \(\Omega(\mathrm{A}_{1})\), and consequently an injective homomorphism of \((\kappa_{\mathrm{A}_{1}} \otimes_{\mathrm{A}} \Omega(\mathrm{A})) / \kappa_{\mathrm{A}_{1}}(1 \otimes df_{1})\) into \(\kappa_{\mathrm{A}_{1}} \otimes_{\mathrm{A}_{1}} \Omega(\mathrm{A}_{1})\). Since \(\kappa_{\mathrm{A}_{1}} \otimes_{\mathrm{A}} \Omega(\mathrm{A})\) is identified with \(\kappa_{\mathrm{A}_{1}} \otimes_{\kappa_{\mathrm{A}}} (\kappa_{\mathrm{A}} \otimes_{\mathrm{A}} \Omega(\mathrm{A}))\), assertion (i'') is therefore equivalent to:

(ii'') the elements \(1 \otimes df_2, \ldots, 1 \otimes df_n\) are linearly independent in \((\kappa_{\mathrm{A}} \otimes_{\mathrm{A}} \Omega(\mathrm{A})) / \kappa_{\mathrm{A}}(1 \otimes df_1)\).

But the conjunction of (ii') and (ii'') is equivalent to (ii), which proves the lemma.

PROPOSITION 6. Let \(k\) be a field of characteristic \(p \neq 0\), \(k'\) be a radical extension of \(k\), of finite degree and height \(\leqslant 1\), and let A be a local regular \(k\)-algebra. Then \(\mathrm{A}_{(k')}\) is a local ring and the following conditions are equivalent:

\begin{enumerate}
\def\labelenumi{(\roman{enumi})}
\item
  the ring \(A_{(k')}\) is regular;
\item
  the \(\kappa_{A}\)-linear map
\end{enumerate}

\[
\kappa_ {\mathrm{A}} \otimes_ {k ^ {\prime p}} \Omega_ {k ^ {p}} (k ^ {\prime p}) \longrightarrow \kappa_ {\mathrm{A}} \otimes_ {\mathrm{A}} \Omega (\mathrm{A})
\]

deduced from the canonical injection \(k^{\prime p}\to \mathrm{A}\) is injective.

Indeed, let \((x_{i})_{i\in \mathrm{I}}\) be a finite \(p\)-basis of \(k'\) over \(k\) (A, V, p.~98); for all \(i\in \mathrm{I}\), set \(f_{i} = x_{i}^{p}\in k\). The \(k\)-algebra \(k'\) is identified with the quotient of \(k[(\mathrm{T}_i)_{i\in \mathrm{I}}]\) by the ideal generated by the polynomials \(\mathrm{T}_i^p - f_i\), hence the A-algebra \(\mathrm{A}_{(k')}\) is the quotient of \(\mathrm{A}[(\mathrm{T}_i)_{i\in \mathrm{I}}]\) by the ideal generated by the polynomials \(\mathrm{T}_i^p - f_i1_{\mathrm{A}}\).

Moreover, \((f_{i})_{i\in\mathrm{I}}\) is a p-basis of \(k^{\prime p}\) over \(k^{p}\), and the \(k^{\prime p}\)-vector space \(\Omega_{k^{p}}(k^{\prime p})\) has the family of \(df_{i}\) as a basis (A, V, p.~97, th. 1). Proposition 6 then follows from Lemma 2.

